\begin{longtable}{|p{\textwidth}|}
\caption{Example CDL description of latlon dataset}
\label{tab:cdl-latlon} \\
\hline \endhead
\hline \endfoot
\hangindent=1.5em\hangafter=1
netcdf latlon example\\
dimensions\\
\hline
\rowcolor{YellowGreen}  time = 1\\
\rowcolor{YellowGreen}  latitude = 360\\
\rowcolor{YellowGreen}  longitude = 720\\
\rowcolor{YellowGreen}  nv = 2\\
\hline

coordinates\\
\hline
\rowcolor{Apricot}\hspace{0.5cm}int32 time (time)\\
\rowcolor{Apricot}\hspace{0.5cm}\hspace{0.5cm}time:axis = "T"\\
\rowcolor{Apricot}\hspace{0.5cm}\hspace{0.5cm}time:bounds = "time\_bnds"\\
\rowcolor{Apricot}\hspace{0.5cm}\hspace{0.5cm}time:coverage\_content\_type = "coordinate"\\
\rowcolor{Apricot}\hspace{0.5cm}\hspace{0.5cm}time:long\_name = "center time of averaging period"\\
\rowcolor{Apricot}\hspace{0.5cm}\hspace{0.5cm}time:standard\_name = "time"\\
\rowcolor{Apricot}\hspace{0.5cm}\hspace{0.5cm}time:units = "hours since 1992-01-01T12:00:00"\\
\rowcolor{Apricot}\hspace{0.5cm}\hspace{0.5cm}time:calendar = "proleptic\_gregorian"\\
\rowcolor{Apricot}\hspace{0.5cm}float32 latitude (latitude)\\
\rowcolor{Apricot}\hspace{0.5cm}\hspace{0.5cm}latitude:axis = "Y"\\
\rowcolor{Apricot}\hspace{0.5cm}\hspace{0.5cm}latitude:bounds = "latitude\_bnds"\\
\rowcolor{Apricot}\hspace{0.5cm}\hspace{0.5cm}latitude:comment = "uniform grid spacing from -89.75 to 89.75 by 0.5"\\
\rowcolor{Apricot}\hspace{0.5cm}\hspace{0.5cm}latitude:coverage\_content\_type = "coordinate"\\
\rowcolor{Apricot}\hspace{0.5cm}\hspace{0.5cm}latitude:long\_name = "latitude at grid cell center"\\
\rowcolor{Apricot}\hspace{0.5cm}\hspace{0.5cm}latitude:standard\_name = "latitude"\\
\rowcolor{Apricot}\hspace{0.5cm}\hspace{0.5cm}latitude:units = "degrees\_north"\\
\rowcolor{Apricot}\hspace{0.5cm}float32 longitude (longitude)\\
\rowcolor{Apricot}\hspace{0.5cm}\hspace{0.5cm}longitude:axis = "X"\\
\rowcolor{Apricot}\hspace{0.5cm}\hspace{0.5cm}longitude:bounds = "longitude\_bnds"\\
\rowcolor{Apricot}\hspace{0.5cm}\hspace{0.5cm}longitude:comment = "uniform grid spacing from -179.75 to 179.75 by 0.5"\\
\rowcolor{Apricot}\hspace{0.5cm}\hspace{0.5cm}longitude:coverage\_content\_type = "coordinate"\\
\rowcolor{Apricot}\hspace{0.5cm}\hspace{0.5cm}longitude:long\_name = "longitude at grid cell center"\\
\rowcolor{Apricot}\hspace{0.5cm}\hspace{0.5cm}longitude:standard\_name = "longitude"\\
\rowcolor{Apricot}\hspace{0.5cm}\hspace{0.5cm}longitude:units = "degrees\_east"\\
\rowcolor{Apricot}\hspace{0.5cm}int32 time\_bnds (time, nv)\\
\rowcolor{Apricot}\hspace{0.5cm}\hspace{0.5cm}time\_bnds:comment = "Start and end times of averaging period."\\
\rowcolor{Apricot}\hspace{0.5cm}\hspace{0.5cm}time\_bnds:coverage\_content\_type = "coordinate"\\
\rowcolor{Apricot}\hspace{0.5cm}\hspace{0.5cm}time\_bnds:long\_name = "time bounds of averaging period"\\
\rowcolor{Apricot}\hspace{0.5cm}float32 latitude\_bnds (latitude, nv)\\
\rowcolor{Apricot}\hspace{0.5cm}\hspace{0.5cm}latitude\_bnds:coverage\_content\_type = "coordinate"\\
\rowcolor{Apricot}\hspace{0.5cm}\hspace{0.5cm}latitude\_bnds:long\_name = "latitude bounds grid cells"\\
\rowcolor{Apricot}\hspace{0.5cm}float32 longitude\_bnds (longitude, nv)\\
\rowcolor{Apricot}\hspace{0.5cm}\hspace{0.5cm}longitude\_bnds:coverage\_content\_type = "coordinate"\\
\rowcolor{Apricot}\hspace{0.5cm}\hspace{0.5cm}longitude\_bnds:long\_name = "longitude bounds grid cells"\\
\hline

data variables\\
\hline
\hspace{0.5cm}float32 OBP (time, latitude, longitude)\\
\hspace{0.5cm}\hspace{0.5cm}OBP:\_FillValue = "9.969209968386869e+36"\\
\hspace{0.5cm}\hspace{0.5cm}OBP:coverage\_content\_type = "modelResult"\\
\hspace{0.5cm}\hspace{0.5cm}OBP:long\_name = "Ocean bottom pressure given as equivalent water thickness"\\
\hspace{0.5cm}\hspace{0.5cm}OBP:units = "m"\\
\hspace{0.5cm}\hspace{0.5cm}OBP:comment = "OBP excludes the contribution from global mean atmospheric pressure and is therefore suitable for comparisons with GRACE data products. OBP is calculated as follows. First, we calculate ocean hydrostatic bottom pressure anomaly, PHIBOT, with PHIBOT = p\_b/rhoConst - gH(t), where p\_b = model ocean hydrostatic bottom pressure, rhoConst = reference density (1029 kg m-3), g is acceleration due to gravity (9.81 m s-2), and H(t) is model depth at time t. Then, OBP = PHIBOT/g + corrections for i) global mean steric sea level changes related to density changes in the Boussinesq volume-conserving model (Greatbatch correction, see sterGloH) and ii) global mean atmospheric pressure variations. Use OBP for comparisons with ocean bottom pressure data products that have been corrected for global mean atmospheric pressure variations. GRACE data typically ARE corrected for global mean atmospheric pressure variations. In contrast, ocean bottom pressure gauge data typically ARE NOT corrected for global mean atmospheric pressure variations."\\
\hspace{0.5cm}\hspace{0.5cm}OBP:coordinates = "time"\\
\hspace{0.5cm}\hspace{0.5cm}OBP:valid\_min = "-2.544442892074585"\\
\hspace{0.5cm}\hspace{0.5cm}OBP:valid\_max = "72.1243667602539"\\
\hspace{0.5cm}float32 OBPGMAP (time, latitude, longitude)\\
\hspace{0.5cm}\hspace{0.5cm}OBPGMAP:\_FillValue = "9.969209968386869e+36"\\
\hspace{0.5cm}\hspace{0.5cm}OBPGMAP:coverage\_content\_type = "modelResult"\\
\hspace{0.5cm}\hspace{0.5cm}OBPGMAP:long\_name = "Ocean bottom pressure given as equivalent water thickness, includes global mean atmospheric pressure"\\
\hspace{0.5cm}\hspace{0.5cm}OBPGMAP:units = "m"\\
\hspace{0.5cm}\hspace{0.5cm}OBPGMAP:comment = "OBPGMAP includes the contribution from global mean atmospheric pressure and is therefore suitable for comparisons with ocean bottom pressure gauge data products. OBPGMAP is calculated as follows. First, we calculate ocean hydrostatic bottom pressure anomaly, PHIBOT, with PHIBOT = p\_b/rhoConst - gH(t), where p\_b = model ocean hydrostatic bottom pressure, rhoConst = reference density (1029 kg m-3), g is acceleration due to gravity (9.81 m s-2), and H(t) is model depth at time t. Then, OBPGMAP= PHIBOT/g + corrections for global mean steric sea level changes related to density changes in the Boussinesq volume-conserving model (Greatbatch correction, see sterGloH). Use OBPGMAP for comparisons with ocean bottom pressure data products that have NOT been corrected for global mean atmospheric pressure variations. GRACE data typically ARE corrected for global mean atmospheric pressure variations. In contrast, ocean bottom pressure gauge data typically ARE NOT corrected for global mean atmospheric pressure variations."\\
\hspace{0.5cm}\hspace{0.5cm}OBPGMAP:coordinates = "time"\\
\hspace{0.5cm}\hspace{0.5cm}OBPGMAP:valid\_min = "7.395928859710693"\\
\hspace{0.5cm}\hspace{0.5cm}OBPGMAP:valid\_max = "82.14805603027344"\\
\hline
\end{longtable}